% =====================================================================
% EJE C - EL DATASET  ---  RESPONSABLE: INTEGRANTE 3
% Que trabajos lo han usado? Que metricas y bajo que protocolo?
% Cual es el mejor resultado publicado? Que problemas conocidos tiene?
% Si el dataset de Kaggle no tiene literatura arbitrada: usar datasets
% ANALOGOS (HKJC, UK, Australia, Japon) y justificar comparabilidad.
% =====================================================================
\subsection{Axis C: The Dataset}
\label{sec:ejeC}

\subsubsection*{\textbf{The dataset specification}} We use the \textit{Horse Racing in HK}
dataset published on Kaggle by Graham Daley, distributed as two tables,
\texttt{races.csv} and \texttt{runs.csv}. We verified the following directly from
the files rather than relying on the dataset page description: 6,349 races and
79,447 individual horse runs, spanning 1997-06-02 to 2005-08-28.
\texttt{races.csv} records per-race attributes (venue, distance, going, class,
sectional and finishing times, dividends); \texttt{runs.csv} records per-horse-per-
race attributes (draw, declared/actual weight, horse rating, gear, trainer and
jockey identifiers, sectional positions, and finishing result). Feature
availability is consistent with a pre-race prediction setting for most columns
(draw, weight, rating, jockey/trainer identity are known before the race), while
several columns (sectional positions, finishing time, dividends) are only known
during or after the race and must be excluded from any predictive feature set.

\subsubsection*{\textbf{Comparable work on Hong Kong racing}} Silverman~\cite{ejeA_gupta2023}
performs a hierarchical Bayesian analysis of Hong Kong horse racing: 36,006
observations from 2,973 races, modelling each horse's running speed and using the
posterior estimates to identify the horse expected to be fastest in each race,
reaching 21.6\% accuracy against the actual outcome. This is, to our knowledge,
the closest jurisdiction-matched published result available to us, same racing
authority (Hong Kong) as our dataset, though roughly half our race count (2,973
versus our 6,349) and with the exact period left unspecified in the source
available to us, so we cannot confirm whether it overlaps with our 1997--2005
window or covers a different span of Hong Kong Jockey Club history. The target
itself also differs subtly from ours: Silverman's model identifies the horse with
the \textit{highest estimated running speed}, which is correlated with but not
identical to the \textit{official race winner} (a horse can run fastest on average
and still lose due to pace, positioning, or trip). We treat this difference as a
documented protocol gap, not a disqualifying one, and use it as our primary
literature baseline. The same author, with Suchard, later moved from this Bayesian
speed-estimation framing toward a discriminative conditional-logit model with
frailty terms, which we read as the same research line converging toward a ranking-
style formulation closer to what we plan to adopt.

\subsubsection*{\textbf{Analogous datasets, for secondary reference}} Beyond the jurisdiction-
matched Silverman result above, we also keep structurally similar horse-racing
datasets from other jurisdictions as secondary analogs, both to triangulate what
accuracy range is achievable with modern classical ML (Silverman's own result
predates tree-ensemble methods) and because the assignment allows analogous
datasets to complement. The closest analogs by scale are Gupta and
Singh~\cite{ejeA_gupta2023,ejeA_gupta2024b} (Turf Club of India, 14,750 records,
2017--2019, 11.8\% positive class) and Davoodi and Khanteymoori~\cite{ejeA_davoodi2010}
(Aqueduct, NY, 100 races). They cover a different racing jurisdiction, wagering
structure, and (in Davoodi and Khanteymoori's case) a much shorter time window, so
any accuracy figure taken from them should be treated as a loose analog, not a
like-for-like target; we discuss this comparability gap explicitly in the baseline.

\subsubsection*{\textbf{Known issues}} Three dataset-level limitations are worth stating plainly.
First, temporal coverage ends in 2005, so any claim about current racing dynamics
(field sizes, betting patterns, track conditions) should be qualified as
historical. Second, the positive-class rate is structurally low (roughly one winner
per race of typically 10--14 runners), mirroring the imbalance documented in
Gupta and Singh~\cite{ejeA_gupta2023} for a different jurisdiction, so imbalance
handling is not optional. Third, several columns in \texttt{runs.csv} (such as
gear, sectional positions for shorter races) contain structurally missing values
which must be documented and imputed (or excluded) accordingly rather than dropped
by default.